% Registered before the corresponding data was gathered; the public commit
% history of this file is the receipt (cited by date in Methodology).
% Quantities measured later bind to the decision rules stated here.
\section{Registered hypotheses}
\label{sec:hypotheses}

Three questions in this work are decision-grade: their answers change what
the consuming inference project builds next. Each is registered here ---
hypothesis, operationalisation, predictions, decision rule --- before the
corresponding data is gathered, with this repository's commit history as
the receipt. The remaining verification gates in the paper are ordinary
methodology checks and are deliberately not elevated to this status.

\subsection{Issue-coupled costs and differential projection}
\label{sec:hyp-coupling}

This hypothesis was formed after the pipe-binding measurements (which
placed \texttt{IDP.4A} on the fma pipe, contending with FFMA issue slots)
and is registered before any differential-projection data exists.

\emph{Projected cost deltas between kernel variants that shift work onto a
contended issue pipe carry larger projection error than deltas that shift
the same work onto an uncontended pipe; under a purely additive cost model
the excess swallows the flash-attention variant-arbitration margin.}

Operationalisation: one synthetic base kernel; graded injections of FFMA
and \texttt{IDP.4A} (fma pipe, contended), \texttt{LOP3} (alu pipe,
uncontended at the operating point), and \texttt{POPC} (separate
quarter-rate unit); injection sizes chosen so true deltas land in the
5--25\% band, the arbitration operating range. Two cost models are
compared: \emph{naive additive} (op count \(\times\) reciprocal
throughput, summed) and \emph{per-pipe max} (per-pipe issue demand,
maximum across pipes, capped at the 4-per-SM-cycle issue rate).

Registered predictions: (i) additive-model differential error is larger
for fma-pipe injections than for alu-pipe injections; (ii) the per-pipe-max
model removes most of that excess; (iii) given (ii), the per-pipe-max
differential error bound discriminates deltas of 10\% and above.

Decision rule: the flash-attention variant arbitration proceeds by
projection only if the measured per-pipe-max differential error bound is
below the projected inter-variant delta; otherwise the projection route is
abandoned for that decision and both variants are implemented and measured.

\subsection{A hand-rolled NVLink exchange against the NCCL call floor}
\label{sec:hyp-exchange}

Registered before any interconnect measurement exists.

\emph{A peer-store-plus-fence exchange primitive over NVLink completes a
two-GPU exchange of at most 20\,KiB in less than half the steady-state
NCCL all-reduce per-call time at the same size.}

Operationalisation: the primitive is a peer \texttt{STG},
\texttt{\_\_threadfence\_system()}, and a system-scope flag poll, measured
end-to-end and litmus-tested for correctness (message passing with a data
check). The comparator is a two-rank NCCL all-reduce over 4--64\,KiB with
\texttt{NCCL\_ALGO}, \texttt{NCCL\_PROTO}, and channel count pinned and
recorded, taken at steady state (cold-channel rows are separate). Before
the comparison is made, a composed prediction of the primitive's round
trip from its constituent table rows (peer-store latency, fence cost, poll
read) must agree with the end-to-end measurement within a stated tolerance
--- the interconnect analogue of the projection gate.

Decision rule: the deterministic peer-exchange mechanism in the consuming
project proceeds only if the validated primitive beats half the freshly
re-measured NCCL per-call floor at 20\,KiB and below; otherwise the
mechanism is recorded as not viable on this fabric.

\subsection{Dispatch-cost composition and the decode ceiling}
\label{sec:hyp-dispatch}

Registered before any launch-family measurement exists.

\emph{The launch-family rows (empty-kernel dispatch, graph-node replay,
event costs) compose to predict the per-token host dispatch overhead of a
production-shaped decode graph within \(\pm20\%\), and the graph-replay
rows predict the decode-rate ceiling that full capture would unlock.}

Operationalisation: per-token dispatch overhead is node count
\(\times\) the relevant per-node row, compared against a profiler baseline
measured freshly on the same host and driver --- prior-session anchors are
re-measured, never reused. The ceiling prediction divides measured compute
time by the sum of compute time and predicted captured dispatch time.

Decision rule: the consuming project's capture work is prioritised by the
predicted ceiling; a composed prediction outside \(\pm20\%\) blocks the
use of these rows for that prioritisation and is published as a model
failure.
